\begin{figure}
\caption{Resultados do desempenho F1 dos circuitos 22 ao 31 com diferentes tamanhos de vetor de entrada na base \textit{ChatGPT} - Fácil. Esses modelos implementam modelos de ensemble através de \textit{hard} ou \textit{soft voting} de circuitos quânticos. As barras comparativas do lado direito exibem a diferença percentual da métrica F1 ao se reduzir de 768 atributos para 16 atributos, divididos pela profundidade de cada circuito. As marcações que indicam semelhança estatística são advindas do teste pareado de Wilcoxon com correção de p-valor de Holm.}
\label{fig:n_features:easy:g7}
     \centering
    \centering
    \resizebox{!}{0.90\textheight}{%  <-- Ajuste a altura aqui
        \import{capitulos/4_resultados/figuras}{n_features_chatgpt_easy_g7.pgf}
    }
      \par\medskip\ABNTEXfontereduzida\selectfont\textbf{Fonte:} Elaborada pelo autor (2025) \par\medskip
\end{figure}